\textbf{Phần 1. Trắc nghiệm lựa chọn}
\setcounterpageref{ex}{0}

\begin{ex} %[2X13N67]
Để kiểm đếm số lượng các bạn trong lớp sinh vào các tháng, ta tiến hành
\choice
{Đo chiều cao}
{\True Thu thập ngày tháng sinh, phân loại theo tháng và đếm}
{Đo cân nặng}
{Làm bài kiểm tra}
\loigiai{
Thu thập ngày sinh, phân loại theo tháng và đếm số lượng.\\
\textbf{Đáp án của câu hỏi là B}
}
\end{ex}

\begin{ex} %[2X13N67]
Trong một giờ học thực hành, nhóm của Lan tung đồng xu 10 lần và ghi lại số lần xuất hiện mặt Sấp và mặt Ngửa. Hoạt động này là
\choice
{Đo độ dài}
{Xếp hình khối}
{\True Thực hành kiểm đếm và ghi chép số liệu}
{Vẽ đoạn thẳng}
\loigiai{
Đó là thực hành kiểm đếm và ghi chép số liệu.\\
\textbf{Đáp án của câu hỏi là C}
}
\end{ex}

\begin{ex} %[2X13H67]
Tung đồng xu 10 lần thu được 6 lần mặt Sấp. Số lần xuất hiện mặt Ngửa là
\choice
{\True $4$ lần}
{$3$ lần}
{$5$ lần}
{$6$ lần}
\loigiai{
$10 - 6 = 4$ lần.\\
\textbf{Đáp án của câu hỏi là A}
}
\end{ex}

\begin{ex} %[2X13H67]
Gieo xúc xắc 20 lần ghi nhận được: Mặt 1 chấm xuất hiện 3 lần, mặt 2 chấm: 4 lần, mặt 3 chấm: 5 lần, mặt 4 chấm: 2 lần, mặt 5 chấm: 3 lần. Số lần xuất hiện mặt 6 chấm là
\choice
{$2$ lần}
{$4$ lần}
{\True $3$ lần}
{$5$ lần}
\loigiai{
Tổng số lần mặt 1 đến 5: $3 + 4 + 5 + 2 + 3 = 17$ lần.\\
Mặt 6 chấm: $20 - 17 = 3$ lần.\\
\textbf{Đáp án của câu hỏi là C}
}
\end{ex}

\begin{ex} %[2X13H67]
Trong bảng số liệu trên, mặt nào xuất hiện nhiều lần nhất?
\choice
{\True Mặt 3 chấm ($5$ lần)}
{Mặt 2 chấm}
{Mặt 1 chấm}
{Mặt 5 chấm}
\loigiai{
Mặt 3 chấm xuất hiện 5 lần (nhiều nhất).\\
\textbf{Đáp án của câu hỏi là A}
}
\end{ex}

\begin{ex} %[2X13H67]
Mặt nào xuất hiện ít lần nhất?
\choice
{Mặt 1 chấm}
{Mặt 6 chấm}
{Mặt 5 chấm}
{\True Mặt 4 chấm ($2$ lần)}
\loigiai{
Mặt 4 chấm xuất hiện 2 lần (ít nhất).\\
\textbf{Đáp án của câu hỏi là D}
}
\end{ex}

\begin{ex} %[2X13V67]
Khảo sát món ăn yêu thích của 30 bạn học sinh: 12 bạn thích gà rán, 10 bạn thích xúc xích, còn lại thích pizza. Số bạn thích pizza là
\choice
{\True $8$ bạn}
{$7$ bạn}
{$9$ bạn}
{$10$ bạn}
\loigiai{
$30 - (12 + 10) = 30 - 22 = 8$ bạn.\\
\textbf{Đáp án của câu hỏi là A}
}
\end{ex}

\begin{ex} %[2X13V67]
Số bạn thích gà rán nhiều hơn số bạn thích pizza là
\choice
{$3$ bạn}
{$5$ bạn}
{$2$ bạn}
{\True $4$ bạn}
\loigiai{
$12 - 8 = 4$ bạn.\\
\textbf{Đáp án của câu hỏi là D}
}
\end{ex}

\begin{ex} %[2X13V67]
Khảo sát phương tiện đi học của lớp 2B (có 32 bạn): 18 bạn đi xe máy (bố mẹ chở), 8 bạn đi bộ, còn lại đi xe đạp. Số bạn đi xe đạp là
\choice
{$5$ bạn}
{\True $6$ bạn}
{$7$ bạn}
{$8$ bạn}
\loigiai{
$32 - (18 + 8) = 32 - 26 = 6$ bạn.\\
\textbf{Đáp án của câu hỏi là B}
}
\end{ex}

\begin{ex} %[2X13V67]
Số bạn đi xe máy nhiều hơn số bạn đi bộ là
\choice
{$8$ bạn}
{$12$ bạn}
{\True $10$ bạn}
{$14$ bạn}
\loigiai{
$18 - 8 = 10$ bạn.\\
\textbf{Đáp án của câu hỏi là C}
}
\end{ex}

\begin{ex} %[2X13V67]
Thực hành quay chong chóng 15 lần: màu đỏ quay được 8 lần, màu xanh quay được 7 lần. Sự kiện xuất hiện nhiều hơn là
\choice
{Màu xanh}
{Hai màu bằng nhau}
{Không xác định được}
{\True Màu đỏ}
\loigiai{
Màu đỏ xuất hiện 8 lần (nhiều hơn).\\
\textbf{Đáp án của câu hỏi là D}
}
\end{ex}

\begin{ex} %[2X13C67]
Trong một trò chơi rút thăm có thưởng, có 20 lá thăm gồm: 5 thăm trúng gấu bông, 7 thăm trúng hộp bút, còn lại là thăm chúc may mắn lần sau. Số thăm chúc may mắn lần sau là
\choice
{$7$ lá}
{\True $8$ lá}
{$9$ lá}
{$10$ lá}
\loigiai{
$20 - (5 + 7) = 20 - 12 = 8$ lá thăm.\\
\textbf{Đáp án của câu hỏi là B}
}
\end{ex}

\textbf{Phần 2. Trắc nghiệm đúng sai}
\setcounterpageref{ex}{0}

\begin{ex} %[2X13H67]
Khảo sát loài hoa yêu thích của 25 bạn: Hoa hồng: 10 bạn, Hoa hướng dương: 8 bạn, Hoa sen: 7 bạn. Xét tính đúng sai:
\choiceTF
{\True Hoa hồng được nhiều bạn yêu thích nhất (10 bạn).}
{\True Hoa sen có ít bạn yêu thích nhất (7 bạn).}
{Hoa hồng nhiều hơn hoa hướng dương 3 bạn.}
{\True Hoa hồng nhiều hơn hoa hướng dương 2 bạn ($10 - 8 = 2$).}
\loigiai{
\begin{itemize}
    \item Hồng nhiều nhất: Đúng.
    \item Sen ít nhất: Đúng.
    \item Hơn 2 bạn (không phải 3). Ý 3 Sai, Ý 4 Đúng.
\end{itemize}
\begin{itemize}
    \item a) \textbf{Đúng}.
    \item b) \textbf{Đúng}.
    \item c) \textbf{Sai}.
    \item d) \textbf{Đúng}.
\end{itemize}
}
\end{ex}

\begin{ex} %[2X13V67]
Gieo đồng xu 20 lần thu được 12 lần mặt Sấp và 8 lần mặt Ngửa. Xét tính đúng sai:
\choiceTF
{\True Mặt Sấp xuất hiện nhiều hơn mặt Ngửa.}
{\True Mặt Sấp nhiều hơn mặt Ngửa 4 lần ($12 - 8 = 4$).}
{Tổng số lần gieo là 24 lần.}
{\True Tổng số lần gieo là $12 + 8 = 20$ lần.}
\loigiai{
\begin{itemize}
    \item Sấp nhiều hơn: Đúng.
    \item Sấp hơn Ngửa 4 lần: Đúng.
    \item Tổng là 20 lần (không phải 24). Ý 3 Sai, Ý 4 Đúng.
\end{itemize}
\begin{itemize}
    \item a) \textbf{Đúng}.
    \item b) \textbf{Đúng}.
    \item c) \textbf{Sai}.
    \item d) \textbf{Đúng}.
\end{itemize}
}
\end{ex}

\textbf{Phần 3. Trả lời ngắn}
\setcounterpageref{ex}{0}

\begin{ex} %[2X13H67]
Khảo sát số con vật yêu thích của 20 bạn: Chó có 9 bạn, Mèo có 7 bạn, Thỏ có 4 bạn. Có bao nhiêu bạn thích Chó?
\shortans{9}
\loigiai{
\textbf{Đáp số:} $9$\\
9 bạn.
}
\end{ex}

\begin{ex} %[2X13H67]
Số bạn thích Mèo nhiều hơn số bạn thích Thỏ bao nhiêu bạn?
\shortans{3}
\loigiai{
\textbf{Đáp số:} $3$\\
$7 - 4 = 3$ bạn.
}
\end{ex}

\begin{ex} %[2X13V67]
Lớp 2A có 35 bạn, trong đó có 18 bạn nam. Hỏi lớp 2A có bao nhiêu bạn nữ?
\shortans{17}
\loigiai{
\textbf{Đáp số:} $17$\\
$35 - 18 = 17$ bạn.
}
\end{ex}

\begin{ex} %[2X13V67]
Tung xúc xắc 15 lần, mặt chẵn xuất hiện 8 lần. Hỏi mặt lẻ xuất hiện bao nhiêu lần?
\shortans{7}
\loigiai{
\textbf{Đáp số:} $7$\\
$15 - 8 = 7$ lần.
}
\end{ex}

\begin{ex} %[2X13V67]
Trong giỏ có 12 quả táo, 8 quả lê và 5 quả xoài. Tính tổng số quả trong giỏ.
\shortans{25}
\loigiai{
\textbf{Đáp số:} $25$\\
$12 + 8 + 5 = 25$ quả.
}
\end{ex}

\begin{ex} %[2X13C67]
Khảo sát 40 học sinh về môn thể thao yêu thích: Bóng đá có 18 bạn, Bơi lội có 12 bạn, Cầu lông có 6 bạn, còn lại thích Bóng rổ. Hỏi có bao nhiêu bạn thích Bóng rổ?
\shortans{4}
\loigiai{
\textbf{Đáp số:} $4$\\
$40 - (18 + 12 + 6) = 40 - 36 = 4$ bạn.
}
\end{ex}
